\Qhead{1.5}{Колебательные системы с двумя и многими степенями свободы}{Колебательные системы с двумя и многими степенями свободы. Нормальные колебания. Вынужденные колебания.}
\Subfile{1.5.1}{sections/01_oscillations/q05/s01}
\Subfile{1.5.2}{sections/01_oscillations/q05/s02}

\recap{Квадратичные энергии задают матрицы $M,K$. Собственная задача определяет частоты и формы, массовая ортогональность --- нормальные координаты. Вынужденное возбуждение определяется проекцией силы; потери ограничивают резонанс.}
\board{$\det(K-\omega^2M)=0$; две формы связанных осцилляторов; расщепление частот; модальный знаменатель.}
\oral{Может ли система двигаться, если наблюдаемая координата обращается в нуль?}
\bait{Можно оставить антирезонанс: в измеряемой координате компенсируются вклады разных мод, а не исчезает вся энергия.}
\sources{Балакший~\cite{balakshy}, рукописный конспект: PDF-с.~19 (номер листа~17), квадратичные энергии и матричная задача. Вопрос~5 коллеги --- альтернативное объяснение через два маятника. Массовая нормировка, биения и вынужденный отклик выведены из указанных уравнений.}
